%% --- THAY khối \title ... \keywords trong sn-article.tex bằng nội dung dưới ---
%% Điền tên, email, \affil theo thực tế; giữ một dòng \documentclass phù hợp (vd. sn-mathphys hoặc sn-aps).

\title[NL2SQL Multi-Agent]{A Multi-Agent LLM Architecture for Natural Language Access to Relational Business Data}

\author*[1]{\fnm{Given} \sur{Surname}}\email{you@institution.edu}

\affil*[1]{\orgdiv{Department}, \orgname{University}, \orgaddress{\city{City}, \country{Country}}}

\abstract{Natural Language to SQL (NL2SQL) can reduce the barrier between business users and relational data, but large language model (LLM) systems still struggle with schema grounding, join selection, aggregation, and nested logic. This paper presents an application-oriented architectural study of a six-step multi-agent NL2SQL pipeline with explicit reasoning checkpoints: Question Analysis, Schema Selection, Query Planning, SQL Generation, SQL Refinement, and SQL Validation. The system is evaluated under a fixed protocol on the Spider 1.0 development split (1,034 questions), which we use because the official Spider 1.0 submission server no longer accepts new submissions. On this full development set, the proposed 6-step pipeline achieves 77.8\% Exact Match (EM) and 85.6\% Execution Accuracy (EX), improving over a matched 4-step baseline with 73.7\% EM and 81.2\% EX. This evidence supports the practical value of reasoning decomposition for robustness and controllability in NL2SQL, while keeping the paper within the scope of a controlled dev-set architectural study rather than a new official leaderboard claim. Mixed-model configuration and reproducibility details are summarized in Table~\ref{tab:repro}.}

\keywords{NL2SQL, Text-to-SQL, multi-agent systems, large language models, business analytics, decision support}
